\section{The General Principle of Good Administration as a Constitutional Safeguard}

General principles of EU law represent `general, comprehensive [...], and naturally inherent'\autocite[42]{audiolux} constitutional principles that permeate the Union legal order.\autocite[7]{tridimasGeneralPrinciplesEU2020} 
They appeared as a natural consequence of the political and legal identity of the European Union,\autocite[2]{TEU} a nascent legal order created by international treaties that provide a framework for incremental integration among the MSs.\autocite[17–18]{tridimasGeneralPrinciplesBook2006}

The 1957 Treaty of Rome already empowered the Court of Justice to review legal acts,\autocite[173]{treatyofRome} but the Treaty did not provide detailed guidance on the standard of review to be applied by the court.\autocite[265–266]{craigFoundations2018} 
Therefore, in order to fill in the gaps, and in the interest of ensuring that Union institutions would face scrutiny,\autocite[269]{craigFoundations2018} the Court of Justice developed a system of background principles of law inferred from the legal and constitutional traditions of the MS.\autocite[269–270]{craigFoundations2018}
The work of the Court in this regard was not a mere comparative study of the principles that underpinned MS' legal systems, but rather an investigation of `best' interpretation of shared values in light of the objectives of the Treaties.\autocites[270]{craigFoundations2018}[See also][point 69]{audioluxAG}

General principles of EU law concern both principles derived from the concept of the rule of law as understood in a liberal democracy,\autocites[9]{tridimasGeneralPrinciplesEU2020}[point 69]{audioluxAG} and principles derived from the institutional design of the EU.\autocite[11–12]{tridimasGeneralPrinciplesEU2020}
This chapter will focus on the general principle of good administration, its history, and its role as a constitutional safeguard in the European composite administrative enforcement model.

\newpage
\subsection{The Concept of 'Good Administration'}

`Good administration' appears in the Union legal order both as a general principle of EU law and as a fundamental right enshrined in the CFR. This section will deal with the two sources of this right in the EU's legal system, and provide an analysis between the two.

\subsubsection{Good Administration as a General Principle of EU Law}

One of the first instances where the Court of Justice undertook the task of `filling in the gaps' left by the Treaties was in \emph{Algera}, one of the earliest cases of the Court, dealing precisely with administrative matters:

\begin{quote}
    The possibility of withdrawing such measures is a problem of administrative law, which is familiar in the case-law and learned writing of all the countries of the Community, but for the solution of which the Treaty does not contain any rules.
    Unless the court is to deny justice it is therefore obliged to solve the problem by reference to the rules acknowledged by the legislation, the learned writing and the case law of the member countries.\autocite[para A III]{algera}
\end{quote}

Throughout the decades, the principle of good administration has been greatly developed by the CJEU.
The principle appears in the case law in various forms: as ``good'',\autocite[para 1A]{alvis} ``sound'',\autocite[para A II]{societesarre} ``proper''\autocite[paras 7 and 12]{burban} administration. 
While today the concept of good administration is widely recognized in the jurisprudence to constitute a general principle of Union law, its early years were marked by a reluctance to recognize it as such.\autocite[66–68]{mihaescu-evansRightGoodAdministration2015} 

The resistance may have stemmed from the fact that the concept of good administration was referred to in inconsistent ways by the judiciary. 
For example, good administration appeared in the case law as an ``obligation'',\autocite[paras 245 and 279]{volkswagenAg} a ``duty'',\autocite[126]{bundesverbandDeutscherBanken} a ``rule'',\autocite[4]{vanEick} a ``requirement'',\autocite[71]{vba} or a ``principle''.\autocite[24–31]{tetraPak} 
A narrow interpretation of the principle is also observed in cases where the Court of Justice held that while violations of the principle of good administration might have occurred, they did not have an effect on the final legality of the act in question.\autocite[67–68]{mihaescu-evansRightGoodAdministration2015}

Nowadays, the principle of good administration is widely recognized as a general principle of Union law.\autocite[210]{hotelCipriani} However, the jurisprudence and the literature concede that good administration is an `umbrella' concept, without a precise delineation as a standalone principle,\autocite[points 60–61]{commVScottAG} and that it is defined by the more precise procedural rights and requirements it encompasses.\autocite[73–74]{mihaescu-evansRightGoodAdministration2015}

To discern the content of the general principle of good administration, we turn to Article 41 CFR, which provides a number of sub-principles that can be understood as falling under the umbrella of good administration.
\autocites[75–76]{mihaescu-evansRightGoodAdministration2015}[127]{Tillack}

\subsubsection{The Charter of Fundamental Rights and Article 41 CFR}

The Charter of Fundamental Rights was first solemnly proclaimed in 2000 as a political document,\autocite[27–28]{fraHandbookDataprotection2018} and given binding force in 2009 with the entry into force of the Lisbon Treaty.\autocite[6(1)]{TEU}
Section V, on citizens' rights, contains Article 41, titled `Right to good administration'.\autocite[41]{CFR}
Article 41 CFR is directly connected to the general principle of good administration established in the case law, and based on the `existence of the Union as subject to the rule of law'.\autocite[Explanation on Article 41]{ExplanationsCFR}

Article 41 CFR provides a number of rights inserted under the umbrella of `good administration'.
As the CJEU outlined in \emph{Tillack}:

\begin{quote}
    [The principle of sound administration] constitutes the expression of specific rights such as the right to have affairs handled impartially, fairly and within a reasonable time, the right to be heard, the right to have access to files, or the obligation to give reasons for decisions, for the purposes of Article 41 of the Charter of fundamental rights of the European Union\autocite[127]{Tillack}
\end{quote}

The wording of the Court hints at a non-exhaustive list of principles,\autocite[91–92]{hofmannEvansRelationCharterFundamental2013} similarly to Article 41(2) CFR, which states that the right `includes' a number of specific procedural rights.\autocite[41(2)]{CFR}

One important note to make is that Article 41(1) CFR explicitly limits the institutional scope of the Charter right to good administration.
The article circumscribes the right to good administration to the `institutions, bodies, offices and agencies of the Union'.\autocite[41(1)]{CFR}
The Charter right also presents a narrower material scope compared to the case–law principle, making references to the `individual', the right to access `his or her' file, or an `individual measure'.\autocite[86–87]{hofmannEvansRelationCharterFundamental2013}

Therefore, attention must be paid to the way these two sources of the right to good administration interact with each other.

\subsubsection{The Coexistence of Good Administration as a General Principle and as a CFR Right}

The relationship between good administration as a general principle of Union law and as a Charter right has been the subject of discussion both in the literature and in the case law.\autocite{mihaescu-evansRightGoodAdministration2015}

As highlighted above, the right to good administration contained in Article 41 CFR is narrower in scope than the general principle of good administration as developed in case law.\autocite[148]{mihaescu-evansRightGoodAdministration2015}
The limitations in the personal, material and institutional scope of the Charter right are problematic especially because, as it is the case with data protection law, MSs are most often the enforcers of Union law.\autocite[150]{mihaescu-evansRightGoodAdministration2015}

The CJEU has however clarified that while the right contained in the Charter does not apply to MSs,\autocite[68–69]{ysVMinister} when invoked as a general principle, good administration must be observed by MSs when acting in the scope of Union law.\autocite[50]{HNvMinisterJustice} 

\subsection{Good Administration in the European Composite Administrative Enforcement Model}

As explored earlier, the EU presents a unique institutional design, where the Union legal order supplements the national legal orders of the MSs. 
Consequently, MSs are often the main enforcers of Union law, which is supposed to be applied uniformly and predictably across the territory of the Union.
In the realm of administrative enforcement, the constitutional structure of the EU generates a ``multi-level'', or ``composite'' administrative procedure, where the administrations of different MSs and Union institutions might be collectively responsible for carrying out enforcement on an individual case.\autocite[91–92]{mihaescu-evansRightGoodAdministration2015} 

This section will highlight the features of the `integrated' administrative model of the EU, and the role of the principle of good administration in ensuring the rule of law across the Union legal order.

\subsubsection{The Composite Administrative Enforcement Model}

EU law presents many policy areas where enforcement is shared among different institutional actors at the Union and national level.\autocite[138]{hofmannCompositeDecisionMaking2009}
We can distinguish between two types of cooperation: horizontal, between actors on the same level of government (such as National Competition Authorities of different MSs), and vertical, between actors on different levels of government (such as a National Competition Authority and the European Commission).

The `composite' nature of administrative procedure in the EU is a natural by-product of a decentralized European administration governed by subsidiarity and transjurisdictional cooperation.\autocite[91–92]{mihaescu-evansRightGoodAdministration2015}
However, issues arise from the convergence of different administrative procedural requirements established by the different legal frameworks the actors are subject to, with effects on individual cases and the overall effectiveness of judicial oversight of administrative action.\autocite[137]{hofmannCompositeDecisionMaking2009} Hofmann observes how the lack of systematic harmonization of administrative law exacerbates this problem in the EU.\autocite[137]{hofmannCompositeDecisionMaking2009}

The principle of good administration however could function as a tool to fill in the gaps left by the lack of comprehensive harmonization, acting as a tool to ensure the respect of the rule of law across the complex administrative procedures taking place in the EU.\autocites[92–93]{mihaescu-evansRightGoodAdministration2015}[351]{nehlGoodAdminProceduralRight2009}

\subsubsection{The `Judicialisation' of the Administration}

The right to good administration is explicitly linked to the right to effective judicial protection established by Article 47 CFR.\autocites[Article 41]{ExplanationsCFR}[79–80]{mihaescu-evansRightGoodAdministration2015}
The principles of good administration mirror the principles of effective judicial protection, which is also viewed as an `umbrella' concept.\autocite[79–80]{mihaescu-evansRightGoodAdministration2015}
Establishing strong procedural safeguards at the administrative stage is crucial to ensure the rule of law is upheld during the administrative procedure, and to enable an appellant to successfully challenge the actions of the administration at the judicial stage.\autocite[79–80]{mihaescu-evansRightGoodAdministration2015}
For example, an appellant's right to effective judicial protection would be undermined if a judge would not be able to ascertain the logic and reasoning behind a decision of the administration, hence the need for the administration to provide reasoned decisions.\autocite[81]{mihaescu-evansRightGoodAdministration2015}

The importance of upholding the principle of good administration as a feature of the administrative system is also augmented by the increasing complexity of the European administrative model.\autocite[86–87]{mihaescu-evansRightGoodAdministration2015}
While it had been accepted by the Courts that the right to effective judicial protection would not be infringed upon as long as an ex-post judicial review of administrative acts was available, the insitutional complexity of the EU has made it increasingly difficult to challenge the discretionary power of national administrations.\autocite[85–86]{mihaescu-evansRightGoodAdministration2015}
In fact, it would be impossible for the EU judiciary to substitute itself to national administrative authorities and replace their power of appraisal, given the delicate institutional balance that exists between Union institutions and the MSs.\autocite[87]{mihaescu-evansRightGoodAdministration2015}

However, in a composite administrative system, with national authorities whose decisions are sometimes liable for the protection of the fundamental rights of over 450 million people, integrating good administration as a design of administrative processes may be a counterweight to the discretionary powers of administrations, leading to better outcomes for the rule of law and the protection of fundamental rights.

\subsubsection{A limitation on the administration's discretionary power}



\subsection{The administration's duty to deliver careful and reasoned decisions}

\subsection{The collective duty of Data Protection Authorities to deliver careful and reasoned decisions}

\todo[inline]{talk about subsidiarity and proximity from the hofmann paper on composite decision making}